\chapter{学习建议}
\section{选课策略与原则}
\SourceParagraph{考虑到大二大三的同学已经经历过多次选课，经验十分丰富，本章主要面向刚入校的新同学，讲解一些选课的基本原则与注意事项。}
\SourceParagraph{最重要的事情放到最前面。作为刚入学的大一新生，在选课之前请务必！务必！务必！ 通读南哪助手新生问答\&指南 2.0（2025 版）（\url{https://www.yuque.com/greatnju/q-a2.0}）中“总则部分”-“教务”-“选课与上课”中的全部内容。这一部分写的非常详实，普适，建议大家在初选前期、中期以及退补选阶段都通读至少一次，读得越多越好。}
\SourceParagraph{掌握了选课的基本原理之后，让我们了解一下在南赫学院选课有哪些需要格外注意的事项。南赫学院的课程类型在大一上可以被分为被指选的必修课（必须要上且自己不能选择班级、更改时间）、自己选择的必修课和选修课。具体来说，大一上的指选必修包括：}
\SourceListItem{1，微积分 I；2，英语读写；3，英语听说；4，大学化学（上面这四门课不能更改老师和时间，每门课的学分都不少，是大一上相对而言比较重要的课程）}
\SourceListItem{5，认识地球系统科学（指定选修相当于必须上课）；6，北欧文化与社会；7，习近平新时代中国特色社会主义思想概论（但是你可以参考课程评价红黑榜更改时间和老师）下图是 2026 级大一上学期的指选课程，仅供参考：}
\begin{figure}[H]
  \centering
  \includegraphics[width=0.92\textwidth,height=0.78\textheight,keepaspectratio]{figures/p25\_02.png}
\end{figure}
\SourceParagraph{而需要自选的必修课有体育课、科光和美育课程、以及通识课和悦读。这几门课都是要你自选并等抽签，如果没中签的话，还得退补选重新抢。选择这类课程的时候也需要参考红黑榜。}
\SourceParagraph{规则：科学之光和美育要在大一一年之内修完，也就是说两学期内必须修满一门科学之光和一门美育；体育课在大一到大三六个学期要修四次，即上四个学期的体育课即可，纯通识课分为人文\&科学（四年修完即可），悦读是四年选择三本（不同类型的）书，三本书合成一个学分。（到时候来个 25 级的说明一下）}
\SourceParagraph{上述内容是基于 23\&24 级的进行编写的，据说 25 级及以后会有所调整，请注意官方通告。}
\SourceParagraph{（科光、美育、红黑榜等名词请参考助手的新生指南）}
\SourceParagraph{一些建议：（希望大家不要太累）}
\SourceParagraph{不同的同学适应大学生活的节奏不同，基础也不同。倘如你发现自己可以相对轻松的应付上述课程，可以考虑转专业选修其他自己感兴趣的课程。推荐在力所能及的情况下，选择一门编程课进行选修。对于零基础的同学，建议去学通修计算机的 Python（这个课难度不小，老师给分比较友好），或者通修 C 语言课程。尽量不要选择 C 语言程序设计基础和智能程序设计(C 语言)这门两课，虽然课程上面写着基础二字，但是对于零基础同学极其不友好，学起来很痛苦。}
\SourceParagraph{如果想要敦实数学基础的话，请慎重对待选修其他数学课的想法。慎修微积分 1（第一层次）和数学分析/高等代数等课程，可以选择微积分 1（第二层次），前几门课程难度不小，实际得分可能较低。选课不慎，会拉低自己的绩点，打击自己学习的信心，单纯为了学习知识而选课需要慎重，更多时候旁听会是一个更具性价比的选择。在 nju，从文院到技术科学实验班，绝大多数的学院与老师对旁听都非常鼓励，如果希望学习更多知识的话是十分推荐的。}
\SourceParagraph{同样地，如果有时间，你可以去旁听线性代数（第一层次）课程，旁听学点线代知识还是很好的，也为大一下学线代做点准备，同样建议不要选通修的线代课。}
\SourceParagraph{去年还有同学选了中文的大学化学这门课，事实上南赫自己的化学课程难度对于大多数同学来说已经有一定的难度，建议尽量不要选择化学化工学院的大学化学 A 以及化学实验基础。这两门课的课程体系不同，难度也不小且给分偏低，现工院的大学化学稍好一些，但还是建议旁听，慎重选课。}
\SourceParagraph{关于是否选择跨专业课程的事情，笔者认为是需要合理安排的，如果读者非常希望转专业的话是可以选择的，因为选择了跨专业课程就会在成绩单上留下课程名，未来无论是跨专业保研/出国/考研的话都是可以看到该课程的。但是正是因为选择了跨专业课程并且通过考试会留下课程名，因此对于只是想学习这方面知识充实自我的读者，或许旁听 or 自学是更好的选择，要不然可能会在成绩单下留下一些不太好看的分数并且拉低自己的成绩。当然笔者认为也没必要那么害怕，大学本来就有一个试错的过程，去尝试一下也没什么大不了的，也就是说，想选课就去选课吧。}
\SourceParagraph{一般来说，选课的总学分在 25 分以内为宜，最多 30 个学分。一学期需要考试的课程在 10 门以内，越少选需要考试的课对你的期末周会越友好。警惕一时选课一时爽，期末周火葬场。}
\SourceParagraph{对于退课这件事，开前前两周可以退课，不会留下任何退课记录；开课三到八周退课会留下记录，但是不会有什么实质上的影响，也不会影响升学这种。关于免修不免考，后面会有叙述。}
\section{绩点（GPA）相关}
\SourceParagraph{nju 的绩点算法非常简单易懂，首先需要明确的是，GPA(Grade Point Average，平均学分绩) 指的是课程关于学分加权的分数，详见下图（截自南京大学成绩单说明页）}
\begin{figure}[H]
  \centering
  \includegraphics[width=0.92\textwidth,height=0.78\textheight,keepaspectratio]{figures/p25\_03.png}
\end{figure}
\SourceParagraph{GPA 可以分为三种不同的 GPA：所有课学分绩、学位课学分绩以及保研学分绩。鉴于南赫最大的一届才刚是大三，还没有保研的学生，因此保研学分绩不再叙述（暂时也没办法叙述，现在保研的话可以主要以学位学分绩作为参考）}
\SourceLabel{主要从所有课学分绩和学位课学分绩来进行说明：}
\SourceParagraph{所有课学分绩：顾名思义，即所有课程的学分绩。主要用于出国出示，是南赫学院的评奖评优主要参考 GPA。}
\SourceParagraph{学位课学分绩：学位课程为去除掉选修课和通识课的所有课程，可以粗略地理解为专业课+英语+体育+思政课。可以用于夏令营、预推免等。}
\SourceParagraph{每个学期结束之后，上过的课就会出成绩，我们可以通过南京大学 APP-我的成绩查询自己的课程成绩，这里面已经包含了所有课学分绩和学位课学分绩。而我们也可以查在自己学院的学位课学分绩排名，方法如下：}
\SourceParagraph{查学位学分绩排名：先登录交换生系统，然后进入这个链接：}
\SourceParagraph{\url{http://elite.nju.edu.cn/exchangesystem/index/create?pid=4}}
\SourceParagraph{（如果这个方法不生效了，就打开交换生系统任意的一个项目点进去，都会有排名）}
\section{辅修}
\SourceParagraph{本段向大家介(quan)绍(tui)辅修以及辅修的技巧。部分选课技巧也适用于跨院系修课等情境。}
\SourceParagraph{辅修有两种，一种是得到辅修学位，一种是得到辅修证明，前者需要修够拿到学位的所有学分，满足学位要求，能够在学信网上显示，后者只需要修满对应专业的三十个学分，且不能在学信网上显示。具体的辅修要求详见辅修专业培养方案。辅修学位可以在毕业的时候加在学位证上，辅修证明则只有一张证明。举两个例子说明：}
\SourceLabel{智能科学与技术辅修：}
\begin{figure}[H]
  \centering
  \includegraphics[width=0.92\textwidth,height=0.78\textheight,keepaspectratio]{figures/p25\_04.png}
\end{figure}
\SourceParagraph{如果想要辅修证明，任选 30 学分课程进行辅修即可，但是前提还要上微积分（第一层次）、微积分 2（第一层次）、线性代数（第一层次），这些是额外的要求 .}
\SourceParagraph{如果想要辅修学位的话，那么表里的所有课程都要上完，并且还要完成辅修学位论文，才算完成，同样要完成额外的数学课程要求。}
\SourceParagraph{对于南赫 er 来说，南赫自己的微积分 1、微积分 2 和线性代数不属于任何层次，也不能靠其他课程进行替代，对于某些专业所以如果你真的想辅修的话，你要上两遍微积分 1、两遍微积分 2、两遍线性代数。}
\SourceParagraph{还有一点，辅修记得认准课程号，课程号一定要是相同的才选，只有相同课程号才会认可。中文名字一样，但课程号不同的课程不要选，选了也不会计入辅修学分。}
\SourceBlankLine[2]
\SourceLabel{化学辅修：}
\begin{figure}[H]
  \centering
  \includegraphics[width=0.92\textwidth,height=0.78\textheight,keepaspectratio]{figures/p25\_05.png}
\end{figure}
\SourceParagraph{化学的辅修就没有数学通修课程的要求，只需要完成 30 学分/修满的要求即可，但是化学专业大二到大四在仙林校区，由于时间冲突，辅修基本无望。}
\SourceLabel{Note:}
\SourceListItem{1，由于呢喃低效的行政效率...一些和苏州校区相关的政策可以说是年年变化。23 级入学时苏州校区的这四个专业因为不在辅修课程列表上，所以不能辅修。但 24 级入学时做了调整，变成了可以辅修技科，25级辅修方案新增了自动化（机器人方向）的辅修}
\SourceLabel{2，已确认数字经济专业不能辅修！数字经济学院确定不接受其他学院学生的辅修申请，但是可以跨专业选课）}
\SourceLabel{3，南赫应该也不能跨专业准出}
\SourceBlankLine[2]
\SourceParagraph{南赫学院的课程安排比较饱满，因此辅修学位需要投入大量精力，尤其是本就比较困难的专业，比如数学、物理、天文等，这可能会导致绩点遭受重创。同时生活在苏州校区意味着大部分的课程没法线下参与，只能自学然后考试，除非是苏州校区的几个学院的专业。并且辅修学位或辅修证明实际上对升学、就业没有关键作用，与为之付出的努力远不成正比，作为加分项可能还不如搞搞竞赛科研、绩点高一些、深入学习专业知识或者跟老师打好关系。}
\SourceParagraph{为缥缈的执念而奋斗，让学习充满自己的生活，却可能没有什么面上的收获，你害怕了吗？即使是这样你也要坚持辅修吗？如果你确定要走辅修这条道阻且长的道路，首先我钦佩你的勇气，然后请往下阅读我的南赫辅修心得。}
\SourceParagraph{首先，决定辅修要趁早！不要等到大二才突然冒出想法！尤其是对基础课程有额外要求的专业（比如数学要求数分高代解析几何并且不算入专业课学分）！大一趁还在南京，能上的课就尽量不要拖到未来补！尤其是一些看起来不需要修但是其实是学位课的奇葩课程（比如数学专业的大物实验）。}
\SourceParagraph{其次，提前规划课程，每个学期修几门课，哪些可以拖，哪些必须及时修，哪些只在春/秋学期开，尽早确定。实际选课的时候注意是否与本院课程冲突，尤其是南赫学院那些只开半个学期的课，如果冲突，记得申请免修不免考，每个学期只有两门可以冲突，也就是说，如果你已经申请了两门免修不免考，还有一门冲突不足 1/3，虽然理论上也能上，但是你无法选这门课！同时也要注意，有些老师不接受免修不免考，注意甄别。如果非常幸运，某门课可以不冲突地选中，也请珍惜这种机会，如果免修不免考名额已经用完，这种课程可以跟老师说明情况，协商一下，让你虽然没有申请免修不免考，却享受免修不免考待遇（即没有考勤）。}
\SourceParagraph{最后，合理规划日常时间，尤其是冲突课程或者到了苏州校区之后。当我们不能听课却想修一门课的时候，我们的目的就不再是从课堂上学到什么了，而是如何高效地解决作业、通过考试，学习是学习，考试是考试，牢记这点。如果还关心自己的绩点的话，请务必合理规划需要考试的科目的学习与复习时间，学会接受“知其然不知其所以然”的状态。}
\SourceParagraph{笔者仅仅修了数分 2 数分 3 高代 2 抽代四门课，便渐感无力，于是放弃辅修，选择自学。但是这段经历本身，以及这段经历中结交的老师、同学，都成为了宝贵的财富。相信抗住了辅修折磨的你，一定会成为更强大的南赫人。}
\SourceBlankLine[2]
\section{如何应对全英文教学}
\SourceParagraph{在南赫，除了思政、体育课和通识课之外，都是全英文教学。因此会有很多新生担心全英文教学的问题。在这里先说明一下，我认为完全不必对全英文教学感到担忧，理由有两点：}
\SourceListItem{1. 按照往年的经验，外教给你们上课的比例不高，大多数老师还是中国老师，这些老师虽然上课讲英文，而且也用着英文 ppt 和书本。但是私下交流都是用中文交流的，如果听不懂，直接用中文问就行，并且助教也都是中国人。}
\SourceListItem{2. 就算是外教，一开始他们也会把语速放得很慢，也会给你们一个适应的过程的，作为老师，他们都很热心的。并且不会的也可以问助教。}
\SourceParagraph{但为了以防万一，（比如社恐）等原因，我在下面给出两个解决方案：}
\SourceListItem{1. 虽然我们上课用的是英文教材，但是 nju 一样会给你发中文教材，如中文微积分和线性代数，这两本教材都是可以参考的，如果英文课本的知识点没有看懂，可以用中文的进行对照。化学没有发中文教材，我在下面化学部分已经给了推荐书目了。}
\SourceListItem{2. 可以多练习一下听力，比如多看看美剧，就当闲着没事干的，你会发现你的听力水平会有巨大的提升，听懂英文课程很容易的。}
\SourceParagraph{对于专业术语问题，不用过于担心，见多了就会了（真的）。事实上，大多数人采取的是方式 1，作业和考试无非就是把中文的内容翻译成英语，再背几个术语就行，所以事实上没什么人会在全英教学上遇到真正的问题。}
\SourceBlankLine[2]
\section{免修不免考}
\subsection{申请流程}
\SourceParagraph{请根据申请流程进行申请，学生申请后务必通过邮件提交《南京大学全日制本科生免修不免考申请表》给任课教师（不提交申请表的，任课教师有权审核不通过），需要先由任课老师审核，再由学生所在院系教务员审核，所有审核环节均在网上办理。结业学生由于不能登录教服平台，如需办理免修不免考，下载填写申请表后，发送至本科生院张老师邮箱 zyjwc@nju.edu.cn。}
\subsection{申请说明}
\SourceParagraph{每学期最多申请两门课免修不免考。需要说明的是，免修不免考不等于选课系统里这门课不再和其他课程冲突了。含免修不免考课程在内，最多冲突两门课。}
\SourceParagraph{网上申请后，找任课教师说明情况和申请理由，由任课教师在网上进行审核。（有的老师会叫你写纸质申请，要在教务网下载专门的表格）}
\SourceParagraph{如果冲突课程的免修不免考审核不通过，会自动删除该课程。}
\SourceParagraph{思政军事类课程、体育课程、实践实验类课程不能申请免修不免考。}
\SourceParagraph{免修不免考申请通过后，仍须交作业并参加所有测验和考试。（因老师而异，可以尝试和老师沟通,所以跨专业选课的免修不免考需谨慎，以防因为一个同学不认识而不知道作业通知等）}
\SourceParagraph{免修不免考的申请时间为开课后两周，请在两周内找任课老师办理。（听说第八周之前都是可以，存疑）}
\SourceParagraph{跨专业选修课不能超过五门}
\subsection{为什么要申请免修不免考？}
\SourceParagraph{你已经学习过相关内容，无需再学。}
\SourceParagraph{你知道老师教的很差劲，想自学。}
\SourceParagraph{课程冲突。}
\SourceParagraph{跨校区上课不方便。}
\SourceParagraph{其他。}
\SourceParagraph{本章节主体内容转载自《技科生存指南》，感谢他们允许进行转载。}
\section{一些学习与科研的经验}
\subsection{提前学习与住宿适应}
这里写给希望提前学习/犹豫是否要提前学习的同学：
\SourceParagraph{提前学习对于理科类的科目，最好的方法永远是早培（丘成桐也这么说！）。提前学习能让你形成一个总的大概的认知。这样即使你英语不是很好，也至少可以知道在讲什么部分，防止被滚雪球。这或许也能让你掌握一个比课程内容更高级或者更底层的视角，比如在学微积分前先看数学分析，在学线代前先看高代。}
\SourceParagraph{专门写给大学前没住过宿舍的同学：如果你们大学前都是走读，并且老师经常只在家长群里发消息，请务必有意识的每天看一下群里的消息，无论你是否喜欢辅导员。同时意识到要与同学一起生活的这个问题，无论采用什么样的社交策略，最终达到一个不耗费太多精力，也不漏掉太多信息/机会的效果。}
\SourceBlankLine[2]
\subsection{提问交流与导师联系}

\SourceLabel{1. 如何提问\&如何与学长学姐交流}
\SourceParagraph{可以参考《别像弱智一样提问》\&《提问的智慧》其实不应该有那么多规矩，正常来说大家都是很友善/热心的，但是为了让双方都对话可以正常愉悦进行，以及提高沟通效率，还是写点东西吧。}
\SourceLabel{主要总结如下：}
\SourceLabel{1. 提前做一些准备：}
\SourceParagraph{搜索已有资料（群文件、官网、经验帖）}
\SourceParagraph{看课程介绍、培养方案、通知文件}
\SourceParagraph{自己先形成一个初步理解}
\SourceParagraph{明确自己真正想问的问题}
\SourceParagraph{主要是，被提问的人不知道提问者想问什么，如果直接上来就是问类似于“学长学姐，这个怎么办”之类的问题会让双方都感到疑惑，别人不知道你遇到了什么，也不知道想解决什么问题。}
\SourceListItem{2. 提问的时候尽量要有礼貌，而不是上来就提出问题，但也不必过于卑微，就是正常说话就可以了}
\SourceListItem{3. 尽量组织出一个完整的内容，并且把问题尽量讲清楚，讲明白，比如可以用背景+目标+已知信息+具体问题这种，倒不是说要把提问的程序复杂化，更不是说学长学姐就有架子这种，但是把问题讲清楚确实有利于解答问题，并且也会让双方更快更好的解决问题}
\SourceParagraph{比如“学长/学姐你好，我希望在大气方向继续深造，我看到其他经验贴里说除了课堂成绩，科研/竞赛等也很重要，但是我不知道应该怎么安排，想问一下：1.大一大二期间有没有必要联系老师提前进组，2.如果想保研的话哪些科研/竞赛经历较为重要，3.有没有可以提前学习的知识，谢谢学长/学姐解答！”}

\SourceParagraph{（上面这个问题是我花了 1 分钟瞎编的，可能也不太正确，也不一定适合每个人的 taste，就是大概展示一个笔者认为还可以的例子吧，如果有啥改进建议欢迎提出）}
\SourceListItem{4. 尽量学会用电脑/平板截图，减少拍屏（因为真的很难看清并且影响观感...）}
\SourceListItem{5. 尽量不要把学长学姐当作搜索引擎，而是当成一个有经验的人，并且不要害怕问问题，大家都是这样走过来的，在查询各种资料后还是没搞定问题，feel free to ask}
\SourceParagraph{如果不是问问题，而是希望和学长学姐正常交流的话，就当是朋友一样交流就行了，我相信大部分学长学姐还是很热心的，并且苏州校区是个孤独的地方，他们希望与更多人交流的（应该吧）}
\SourceBlankLine[2]
\SourceLabel{2. 如何找导师/提前进组}
\SourceParagraph{相信很多同学都听说过，本科阶段可以提前进入导师课题组参与科研，也可能希望通过科研经历了解专业方向、提升自己的能力。但是，很多同学会面临几个问题：什么时候应该联系导师？应该选择什么样的导师？如何联系导师？邮件中应该写些什么？}
\SourceParagraph{下面结合个人经验简单介绍一下：}
\SourceListItem{一、什么时候适合联系导师？}
\SourceParagraph{实际上，联系导师并没有一个严格固定的时间节点，具体情况需要结合个人规划和学习状态决定。}
\SourceParagraph{对于大一同学来说，个人认为不必过早焦虑科研问题。大一阶段课程相对较多，更重要的是适应大学生活、打好专业基础，同时也可以充分体验南京的校园生活和城市环境（其实就是好好玩）。除此之外，很多老师主要在苏州校区，大一去的话不太方便。}
\SourceParagraph{一般来说，大二上学期开始尝试联系导师是一个比较合适的时间点。相比大一，大二阶段专业课程逐渐增加，但整体时间安排会更加灵活，同时也有一些类似大创的科研启蒙项目，可以帮助同学们逐渐了解科研过程。因此，如果希望提前接触科研，大二阶段可以开始主动了解导师并尝试联系。}
\SourceParagraph{当然，具体时间仍然需要根据个人情况决定。如果已经明确未来想从事某个方向，并且具备一定兴趣和基础，也可以适当提前。}
\SourceListItem{二、如何选择导师？}
\SourceParagraph{选择导师最直接的方法是查看学院官网，通常学院官网都会提供教师信息，包括：教师姓名、研究方向、个人主页、代表性论文、联系邮箱等。}
\SourceParagraph{同学们可以根据自己的兴趣选择感兴趣的研究方向，然后进一步了解相关老师的研究内容。}
\SourceParagraph{建议不要只看老师的研究方向简介，也可以阅读老师主页上的代表论文，对研究内容形成初步认识。这样在后续联系导师时，也能够体现自己确实了解过老师的研究方向，而不是盲目发送邮件。}
\SourceListItem{三、如何联系导师？}
\SourceParagraph{目前最常见、也是最正式的方式是通过邮件联系导师。除了邮件之外，也可以通过课程/学术会议与老师建立联系：}
\SourceParagraph{不过，对于大多数本科生来说，邮件仍然是最普遍、最正式的方式。}
\SourceListItem{四、邮件中应该写些什么？}
\SourceParagraph{很多同学会担心老师会不会拒绝本科生进入课题组}
\SourceParagraph{这种担心其实很正常。但实际上，很多导师都经历过学生阶段，也理解本科生希望提前接触科研的需求。尤其现在很多年轻导师，他们通常都比较愿意与本科生交流。}
\SourceParagraph{当然，被拒绝也是存在可能的，比如：老师目前没有本科生培养计划；课题组任务较多，没有时间指导；研究方向与个人兴趣不匹配等。}
\SourceParagraph{这些情况都很正常，并不代表个人能力不足。如果老师明确拒绝，可以继续寻找其他合适的导师。}
\SourceLabel{写邮件时，需要注意以下几点：}
\SourceListItem{1. 保持礼貌和简洁；}
\SourceListItem{2. 明确表达希望进入课题组学习的想法；}
\SourceListItem{3. 说明自己为什么对老师的研究方向感兴趣；}
\SourceListItem{4. 简单介绍自己的专业背景和相关经历；}
\SourceListItem{5. 最好附上一份个人简历。}
\SourceParagraph{需要注意的是，不建议一次性群发大量邮件。更好的方式是根据自己的兴趣选择合适的导师，一个一个认真联系。这样不仅更加尊重老师，也能够提高交流效果。}
\SourceParagraph{最后想分享一点个人感受：}
\SourceParagraph{在本科阶段接触科研时，不需要因为“自己什么都不会”而害怕。实际上，本科生进入课题组的主要目的就是学习科研方法、了解研究过程，而不是一开始就能够独立完成复杂工作。}
\SourceParagraph{大多数导师都理解本科生处于学习阶段，一般都会给予悉心的指导（一般都是硕博的学长学姐来指导，但是不用害怕，他们人都很好的）}
\SourceParagraph{如果你对某个研究方向感兴趣，就大胆尝试联系导师。即使第一次没有成功，也可以继续寻找其他机会。科研经历本身就是一个不断尝试和探索的过程，因此不要害怕，大胆去找。}
\SourceBlankLine[2]
\subsection{社交、课程与综合发展}
\SourceParagraph{在大学生活的初始阶段，尤其应当重视并主动拓展社交网络，积极结识新朋友。多数同学均处于重新建立人际关系的阶段，彼此间对社交互动均有较高需求，因而更易建立联系。就专业学习而言，同专业同学之间的基础知识与能力储备往往相差不大，广泛结识同学不仅能够拓宽信息渠道，亦能显著增加参与各类项目、竞赛及实践活动的机会。集体协作有助于共同应对课业难题，联合参与学术活动或社团事务，既有利于综合能力的锻炼，也可提升个人在群体中的可见度与展示平台。}
\SourceParagraph{关于课程学习，课堂讲授仅为知识获取的一个环节，尤其在南赫学院全英文教学环境下，单次授课内容量大且语言门槛较高，绝大多数学生难以在课上实现完全理解。为此，必须系统性地完成课堂笔记记录，并借助人工智能工具进行深入分析；建议同时使用多款 AI 产品进行交叉验证，因单一模型可能存在解释偏差，而多模型间的回答分歧有助于发现潜在错误，避免被误导，从而提升自学效率与（本人就是 DeepSeek chatgpt 豆包 元宝同时使用 虽然晕是真的晕 但准确性。是理解真的透彻）}
\SourceParagraph{更为关键的是，应避免将全部精力局限于课程知识本身。大学与高中教育模式存在本质差异，后者以应试导向为主，而大学更强调自主性、探索性与多元发展。若过度沉浸于课业学习，而忽视课外活动、科研训练及实践环节，则可能导致能力结构失衡。科研参与、学术讲座、创新项目等活动不仅有助于深化专业认知，更能培养批判性思维、团队协作与问题解决能力，这些素养与课程绩点同等重要，甚至在一定程度上决定了未来深造或就业的竞争力。因此，应当在学业成绩与综合能力提升之间寻求动态平衡，合理分配时间与精力，方能充分把握大学阶段的成长机遇。}
\SourceBlankLine[2]
\subsection{竞赛经验}
\SourceLabel{一、一些竞赛}
\SourceParagraph{在大学生活中也会有各种各样的竞赛，同学们无需对于这些竞赛感到眼花缭乱，也无需对这些竞赛有着门槛过高的担忧，如果是为了保研加分的话，有如下几个比赛值得推荐}
\SourceLabel{1. 中国“互联网+”大学生创新创业大赛}
\SourceLabel{2. “挑战杯”全国大学生系列科技学术竞赛}
\SourceLabel{3. 全国大学生数学建模竞赛}
\SourceLabel{4. 美国大学生数学建模大赛（MCM/ICM）}
\SourceParagraph{但是笔者还是想说，如果是为了保研加分，参加比赛的性价比比不上好好学好课内的知识。参加比赛的意义，是让你有机会见识到更广阔的天地以及真正的学到一些你喜欢的知识。所以在参加一个比赛之前，一定要做好足够的思想准备，尽量不要中途而废。}
\SourceBlankLine[2]
\SourceLabel{二、如何平衡课内和竞赛}
\SourceParagraph{客观上说，参加竞赛一定会分散课内学习的精力，如何平衡时间和精力是一个大问题。具体情况因人而异，笔者在这里分享一些自己的个人经验，仅供参考。}
\SourceLabel{1. 善于利用应试技巧}
\SourceParagraph{大学课程（尤其专业课）的成绩有时考察的并不是对于知识点的掌握程度，而是对期中期末试卷的掌握程度。所以就需要考验各位在有限的往年卷和作业题中发挥无限的智慧去应付考试了。大家各显神通来到了南赫，也相信大家有足够的应试技巧去应对大学的考试。}
\SourceLabel{2. 合理安排时间}
\SourceParagraph{具体情况因人而异，笔者认为，在期末考试前一个月应该心中大概有复习的时间规划。每门科目根据自己掌握的情况在考试前一到两周开始复习（速通），可以拿到一个说得过去的成绩。当然若是要追求极致的话，那还是功夫在平时了。}
\SourceLabel{3. 善于借助外物}
\SourceParagraph{君子性非异也，善假于物也。无论是课程还是竞赛，AI，老师，学长等等都是能让你事半功倍的工具。勤加沟通，多加利用，往往是提高效率，平衡课内外的关键。}
\SourceBlankLine[2]
\SourceLabel{三、如何参加竞赛}
\SourceParagraph{参加竞赛的第一步是去了解这个竞赛，主要可以通过百团大战，和老师和学长的沟通还有自己搜索资料等方式去了解。主要应当去了解这个竞赛是干什么的，它有哪些方向和赛道，主要涉及哪些方面的知识。}
\SourceParagraph{第二部分是去尝试联系老师或者参加过这些比赛的学长，去试着参与其中，去学习，参与。即使一开始实力不够或是无法参与核心，但还是可以多去交流，多去学习，由长时间的量变积攒的质变往往在很短时间内完成。}
\SourceParagraph{至于决定参加比赛之后的一些具体细节，还是得具体情况具体分析。但是，我只能说在那之前以及在其中，要多想，要多问。如果你心存怀疑，不妨先做下去，试一试，也许会有意料之外的转机出现。}
\SourceParagraph{笔者非常幸运，在百团大战时偶然了解了学校的 RM，开始学习和探索，机缘巧合之下进入了这支队伍。可能也成为了改变笔者一生的转折点。}
\SourceParagraph{“人生就像巧克力，你永远不知道下一块是什么味道。”那些改变人生的机会往往在计划之外突然冒出。有时也要去关注一下计划之外的机会。百忙之中，下一步闲棋是很必要的。百团大战上的一会驻足，海报前的惊鸿一瞥，聊天时的随口一提，这些机缘巧合往往是一段伟大征途的开始。}
\SourceLabel{四、一些后话}
\SourceParagraph{世上有不知道多少事情，意气风发的人没办成，苦大仇深的人没办成，腰缠万贯的人也没办成，最后叫谁办成了呢，叫一个茶余饭后，不知道哪根筋搭错，说干就干的人给办成了。}
\SourceParagraph{世上阻碍人看见的最大壁障，既不是才情，也不是困境，更不是经费，而是那份上路的理由。这份理由花什么都买不来，然而一切的一切，都只会在那个理由出现之后，以各种莫名其妙的方式出现。}
\SourceParagraph{所以少年，当一个上路的理由摆在你面前时，勇敢去做吧。}
\SourceParagraph{当你决定出发，世界上的某些东西，才变成了路，变成了桥，变成了学费，变成了学校。正是因为你在上路之前，没有怀疑这些东西会不会出现，空着手就那么走了，这些东西，才会以一种其他人根本认不出来的面貌出现。但凡你觉得你还得装上些什么，先攒个什么东西，这些必要的帮助才能被我交易过来。那么这些东西就往往永远都不会出现了。}
\SourceParagraph{祝好运。}
